\chapter{Bee and Wasp Stings}

\section*{Definition and Overview}
Bee and wasp stings are injuries caused when these insects pierce the skin with their stinger and inject venom. For most people a sting causes sharp pain, redness, and swelling that settles within a day or two. For a small minority, stings can trigger a severe allergic reaction called anaphylaxis, which is a medical emergency. In Australia, native and European honeybees, paper wasps, and the introduced European wasp are the main stinging insects, and deaths from stings are rare but do occur, almost always from anaphylaxis.

\section*{Epidemiology}
Stings are common, especially in late summer and autumn, and in outdoor settings, gardens, and around food and rubbish bins. Beekeepers, gardeners, outdoor workers, and children are more often stung. About 1 in 50 to 1 in 100 Australians has a significant allergic reaction to insect stings, and around 1 in 200 has a severe reaction. Severe reactions can occur at any age but are more dangerous in older adults and in people with heart or lung disease. Anaphylaxis from stings is uncommon but causes around 1--2 deaths per year in Australia.

\section*{Aetiology and Pathophysiology}
\begin{itemize}
    \item \textbf{Mechanism:} the insect injects venom through a stinger into the skin; a honeybee leaves its barbed stinger behind and then dies, whereas wasps can sting repeatedly
    \item \textbf{Venom contents:} the venom contains proteins and enzymes that cause pain, inflammation, and local tissue damage
    \item \textbf{Allergic reaction:} in sensitised people, the immune system produces antibodies (IgE) against the venom; a further sting triggers release of histamine and other chemicals, causing hives, swelling, and, in severe cases, anaphylaxis
    \item \textbf{Toxic reaction:} massive numbers of stings (for example, disturbing a large hive) can cause a toxic reaction even without allergy
    \item \textbf{Risk of anaphylaxis:} more likely in people who have previously had a systemic reaction to a sting, and the risk of a severe next reaction is highest when the previous reaction was severe
\end{itemize}

\section*{Clinical Features}
\begin{itemize}
    \item Immediate sharp pain and burning at the sting site
    \item Redness, swelling, and a raised pale area (wheal) that may persist for hours
    \item Itching around the sting
    \item A visible stinger may remain in the skin after a bee sting
    \item Large local reaction: extensive swelling and redness spreading over many centimetres, peaking at 24--48 hours
    \item Systemic allergic reaction (anaphylaxis): hives away from the sting site, facial or throat swelling, difficulty breathing, wheezing, vomiting, dizziness, collapse
\end{itemize}

\section*{Differential Diagnosis}
\begin{itemize}
    \item Other insect bites and stings, such as ants, mosquitoes, spiders, and ticks
    \item Large local reactions to stings, which are not anaphylaxis but cause marked swelling
    \item Cellulitis or bacterial skin infection, which may follow a sting
    \item Contact dermatitis or plant irritation (for example, nettles)
    \item Other causes of sudden breathing difficulty or collapse, such as asthma, food allergy, or a heart event, in people without a clear sting history
\end{itemize}

\section*{When to Seek Medical Care}
Most stings can be managed at home. Seek medical care for a large local reaction, a sting inside the mouth, multiple stings, or if the wound shows signs of infection (spreading redness, pain, discharge, or fever).

\textbf{Call 000 (or local emergency services) immediately for signs of anaphylaxis:}
\begin{itemize}
    \item Difficulty breathing, wheezing, or throat tightness
    \item Swelling of the tongue, face, or lips
    \item Hives or itchy rash away from the sting site
    \item Vomiting, abdominal pain, or diarrhoea
    \item Dizziness, confusion, collapse, or loss of consciousness
\end{itemize}

\section*{Diagnosis and Investigations}
\begin{itemize}
    \item \textbf{History:} the type of insect if known, the site and number of stings, past reactions to stings, and any symptoms beyond the local area
    \item \textbf{Examination:} assessment of the sting site and of the airway, breathing, and circulation in anyone with systemic symptoms
    \item \textbf{Allergy testing:} skin prick tests or blood tests for venom-specific IgE antibodies may be offered after a systemic reaction to confirm the allergy and identify the insect
    \item \textbf{No routine investigations:} local stings need no tests
\end{itemize}

\section*{Management}
\begin{itemize}
    \item \textbf{First aid:} remove a bee stinger quickly by scraping with a fingernail or a card -- do not squeeze it; wasps leave no stinger
    \item \textbf{Cold compress:} apply ice or a cold pack to the sting to reduce pain and swelling
    \item \textbf{Analgesics and anti-itch treatment:} paracetamol or ibuprofen for pain; an oral antihistamine or a corticosteroid cream for itching
    \item \textbf{Elevation:} raise a stung limb to reduce swelling
    \item \textbf{Anaphylaxis treatment:} adrenaline given by auto-injector is the first-line treatment and must be given promptly; call 000 and lie the person flat, raising the legs
    \item \textbf{Hospital care:} emergency treatment for anaphylaxis includes adrenaline, oxygen, fluids, and observation, usually for several hours
    \item \textbf{Referral:} people who have had anaphylaxis or a significant systemic reaction should be referred to an allergy specialist; venom immunotherapy (desensitisation) can markedly reduce the risk of future severe reactions
\end{itemize}

\section*{Complications}
\begin{itemize}
    \item Anaphylaxis -- a life-threatening allergic reaction requiring emergency treatment
    \item Secondary bacterial infection of the sting site (cellulitis, abscess)
    \item Large local swelling that is painful and prolonged
    \item Toxic reaction from many simultaneous stings, causing fever, nausea, muscle breakdown, or kidney problems
    \item A sting inside the mouth or throat can cause dangerous airway swelling
    \item A retained stinger may cause a persistent lump or infection
\end{itemize}

\section*{Prognosis and Outcomes}
For the great majority of people, a sting causes short-lived pain and swelling that settles within 24--48 hours without treatment. Large local reactions can last up to a week but are not dangerous. People who have had anaphylaxis are at risk of recurrence with future stings, but venom immunotherapy reduces the risk of a severe reaction to around 5 out of 100 over five years. With prompt treatment, even severe anaphylaxis has a good outcome; death from a sting is very rare.

\section*{Prevention and Self-Care}
\begin{itemize}
    \item Avoid disturbing hives, nests, and wasp traps; keep food and rubbish covered outdoors
    \item Wear shoes outdoors, especially in grass or near flowering plants
    \item Use insect repellent around outdoor areas and avoid wearing bright clothing and strong perfumes, which attract insects
    \item Keep windows and doors screened, and check for nests in roofs, trees, and under eaves
    \item If a bee or wasp approaches, stay calm and still or move away slowly; do not swat it
    \item People with known venom allergy should carry an adrenaline auto-injector at all times and wear a medical alert bracelet
    \item Ensure family, friends, and school or workplace know how to use the auto-injector and when to call 000
\end{itemize}

\section*{Sources and Further Reading}
The following sources provide reliable, up-to-date information on bee and wasp stings.
\begin{itemize}
    \item NHS. \textit{Insect bites and stings.} https://www.nhs.uk/conditions/
    \item Mayo Clinic. \textit{Bee stings.} https://www.mayoclinic.org
    \item MedlinePlus. \textit{Bee and wasp stings.} https://medlineplus.gov
    \item MSD Manual. \textit{Anaphylaxis.} https://www.msdmanuals.com
    \item Centers for Disease Control and Prevention. \textit{Venomous animals.} https://www.cdc.gov
\end{itemize}
